\section*{ID10131: Celestial Gyromagnetic Motion: γ}
\paragraph{Metadata.}
PiID: 9752148293366 | RTEID: RTE:P39340.9229:T1.9899 | UUID: f1dc9323-4894-5d4a-9afa-306526fcc75f

\paragraph{Canonical Statement.}
Orbital systems exhibit universal gyromagnetic ratio \$\textbackslash{}gamma = L/\textbackslash{}mu\$ where \$L\$ is angular momentum and \$\textbackslash{}mu\$ is magnetic moment from current loop: \textbackslash{}begin\{equation\} \textbackslash{}gamma = \textbackslash{}frac\{L\}\{\textbackslash{}mu\} = \textbackslash{}frac\{mvr\}\{qvr/2\} = \textbackslash{}frac\{2m\}\{q\} \textbackslash{}end\{equation\}

\paragraph{Canonical Equation.}
\[
\gamma = L/\mu
\]

\paragraph{Dictionary.}
Orbital systems exhibit universal gyromagnetic ratio γ = L/μ where L is angular momentum and μ is magnetic moment from current loop [ID10131]

\paragraph{Encyclopedia.}
Celestial Gyromagnetic Motion: γ is an algebraic definition, written γ = L/μ. It is an algebraic definition expressing the quantity directly in terms of the variables on its right-hand side. It is built from μ, L, defining γ. Within the Trawin Zero Point registry it serves as one element of the framework's mathematical narrative.
